\section*{5.1. \ The Contraction Mapping Principle}
\source{gilbarg-trudinger:page-0086}

\begin{definition}\label{gt.def5.1.contraction-mapping}\dcref{contraction mapping}\group{chapter-5}\source{gilbarg-trudinger:page-0086}
A mapping $T$ from a normed linear space $\V$ into itself is called a contraction mapping if there exists a number $\theta < 1$ such that
\begin{equation}\label{gt.eq5.1}\dcref{(5.1)}\group{chapter-5}\source{gilbarg-trudinger:page-0086}
\|Tx - Ty\| \leq \theta \|x - y\| \qquad \text{for all } x, y \in \V.
\end{equation}
\end{definition}

\begin{theorem}\label{gt.thm5.1}\dcref{Theorem 5.1}\group{chapter-5}\source{gilbarg-trudinger:page-0086}
A contraction mapping $T$ in a Banach space $\B$ has a unique fixed point, that is there exists a unique solution $x \in \B$ of the equation $Tx = x$.
\end{theorem}

\begin{proof}\source{gilbarg-trudinger:page-0086}
(Method of successive approximations.) Let $x_0 \in \B$ and define a sequence $\{x_n\}\subset\B$ by $x_n = T^n x_0$, $n = 1, 2, \ldots$. Then if $n \geq m$, we have
\[
\|x_n - x_m\| \leq \sum_{j=m+1}^{n} \|x_j - x_{j-1}\| \qquad \text{by the triangle inequality}
\]
\[
= \sum_{j=m+1}^{n} \|T^{j-1}x_1 - T^{j-1}x_0\|
\]
\[
\leq \sum_{j=m+1}^{n} \theta^{j-1}\|x_1 - x_0\| \qquad \text{by \eqref{gt.eq5.1}}
\]
\[
\leq \frac{\|x_1 - x_0\|\theta^m}{1-\theta} \to 0 \qquad \text{as } m \to \infty.
\]
Consequently $\{x_n\}$ is a Cauchy sequence and, since $\B$ is complete, converges to an element $x \in \B$. Clearly $T$ is also a continuous mapping and hence we have
\[
Tx = \lim Tx_n = \lim x_{n+1} = x
\]
so that $x$ is a fixed point of $T$. The uniqueness of $x$ follows immediately from \eqref{gt.eq5.1}. $\square$
\end{proof}

\begin{remark}\source{gilbarg-trudinger:page-0086}
In the statement of Theorem 5.1, the space $\B$ can obviously be replaced by any closed subset.
\end{remark}

\section*{5.2. \ The Method of Continuity}
\source{gilbarg-trudinger:page-0086,page-0087}

\begin{definition}\label{gt.def5.2.bounded-linear-operator}\dcref{bounded linear mapping}\group{chapter-5}\source{gilbarg-trudinger:page-0086,page-0087}
Let $\V_1$ and $\V_2$ be normed linear spaces. A linear mapping $T:\V_1 \to \V_2$ is bounded if the quantity
\begin{equation}\label{gt.eq5.2}\dcref{(5.2)}\group{chapter-5}\source{gilbarg-trudinger:page-0087}
\|T\| = \sup_{x \in \V_1,\, x \neq 0} \frac{\|Tx\|_{\V_2}}{\|x\|_{\V_1}}
\end{equation}
is finite. It is easy to show that a linear mapping $T$ is bounded if and only if it is continuous.
\end{definition}

\begin{theorem}\label{gt.thm5.2}\dcref{Theorem 5.2}\group{chapter-5}\source{gilbarg-trudinger:page-0087}
Let $\B$ be a Banach space, $\V$ a normed linear space and let $L_0$, $L_1$ be bounded linear operators from $\B$ into $\V$. For each $t \in [0,1]$, set
\[
L_t = (1-t)L_0 + tL_1
\]
and suppose that there is a constant $C$ such that
\begin{equation}\label{gt.eq5.3}\dcref{(5.3)}\group{chapter-5}\source{gilbarg-trudinger:page-0087}
\|x\|_{\B} \leq C\|L_t x\|_{\V}
\end{equation}
for $t \in [0,1]$. Then $L_1$ maps $\B$ onto $\V$ if and only if $L_0$ maps $\B$ onto $\V$.
\end{theorem}

\begin{proof}\source{gilbarg-trudinger:page-0087}
Suppose that $L_s$ is onto for some $s \in [0,1]$. By \eqref{gt.eq5.3}, $L_s$ is one-to-one and hence the inverse mapping $L_s^{-1} : \V \to \B$ exists. For $t \in [0,1]$ and $y \in \V$, the equation $L_t x = y$ is equivalent to the equation
\[
L_s(x) = y + (L_s - L_t)x
\]
\[
= y + (t-s)L_0x - (t-s)L_1x
\]
which in turn, is equivalent to the equation
\[
x = L_s^{-1}y + (t-s)L_s^{-1}(L_0 - L_1)x
\]
The mapping $T$ from $\B$ into itself given by $Tx = L_s^{-1}y + (t-s)L_s^{-1}(L_0 - L_1)x$ is clearly a contraction mapping if
\[
|s-t| < \delta = [C(\|L_0\| + \|L_1\|)]^{-1}
\]
and hence the mapping $L_t$ is onto for all $t \in [0,1]$ satisfying $|s-t| < \delta$. By dividing the interval $[0,1]$ into subintervals of length less than $\delta$, we see that the mapping $L_t$ is onto for all $t \in [0,1]$ provided it is onto for any fixed $t \in [0,1]$, in particular for $t=0$ or $t=1$. $\square$
\end{proof}

\section*{5.3. \ The Fredholm Alternative}
\source{gilbarg-trudinger:page-0087,page-0088,page-0089,page-0090,page-0091}

\begin{definition}\label{gt.def5.3.compact-mapping}\dcref{compact}\group{chapter-5}\source{gilbarg-trudinger:page-0087,page-0088}
Let $\V_1$ and $\V_2$ be normed linear spaces. A mapping $T : \V_1 \to \V_2$ is called compact (or completely continuous) if $T$ maps bounded sets in $\V_1$ into relatively compact sets in $\V_2$ or equivalently $T$ maps bounded sequences in $\V_1$ into sequences in $\V_2$ which contain convergent subsequences. It follows that a compact linear mapping is also continuous but the converse is not true in general unless $\V_2$ is finite dimensional.
\end{definition}

\begin{theorem}\label{gt.thm5.3}\dcref{Theorem 5.3}\group{chapter-5}\source{gilbarg-trudinger:page-0088,page-0089,page-0090}
Let $T$ be a compact linear mapping of a normed linear space $\V$ into itself. Then either (i) the homogeneous equation
\[
 x - Tx = 0
\]
has a nontrivial solution $x \in \V$ or (ii) for each $y \in \V$ the equation
\[
 x - Tx = y
\]
has a uniquely determined solution $x \in \V$. Furthermore, in case (ii), the operator $(I-T)^{-1}$ whose existence is asserted there is also bounded.
\end{theorem}

\begin{lemma}\label{gt.lem5.4}\dcref{Lemma 5.4}\group{chapter-5}\source{gilbarg-trudinger:page-0088}
Let $\V$ be a normed linear space and $\M$ a proper closed subspace of $\V$. Then for any $\theta < 1$, there exists an element $x_\theta \in \V$ satisfying $\|x_\theta\| = 1$ and $\dist(x_\theta, \M) \geq \theta$.
\end{lemma}

\begin{proof}\source{gilbarg-trudinger:page-0088}
Let $x \in \V - \M$. Since $\M$ is closed, we have
\[
\dist(x, \M) = \inf_{y \in \M} \|x-y\| = d > 0.
\]
Consequently there exists an element $y_\theta \in \M$ such that
\[
\|x-y_\theta\| \leq \frac{d}{\theta},
\]
so that, defining
\[
 x_\theta = \frac{x-y_\theta}{\|x-y_\theta\|},
\]
we have $\|x_\theta\| = 1$ and for any $y \in \M$,
\[
\|x_\theta-y\| = \frac{\|x-y_\theta-\|y_\theta-x\|\,y\|}{\|y_\theta-x\|}
\]
\[
\geq \frac{d}{\|y_\theta-x\|} \geq \theta.
\]
The lemma is thus proved. $\square$
\end{proof}

\begin{remark}\source{gilbarg-trudinger:page-0089}
If $\V=\mathbb R^n$, it is clear that one can take $\theta=1$ by choosing $x_0$ orthogonal to $\M$.
This will also be possible in any Hilbert space but in general Lemma 5.4, which asserts the existence of a ``nearly orthogonal'' element to $\M$, cannot be improved to allow $\theta=1$.
\end{remark}

\begin{proof}\source{gilbarg-trudinger:page-0089,page-0090}
\textit{Proof of Theorem 5.3.} \quad It is convenient to split our proof into four stages.

\begin{enumerate}
\item[(1)] \textit{Let $S=I-T$ where $I$ is the identity mapping and let $\mathcal N=S^{-1}(0)=\{x\in\V\mid Sx=0\}$ be the null space of $S$. Then there exists a constant $K$ such that}
\begin{equation}\label{gt.eq5.4}\dcref{(5.4)}\group{chapter-5}\source{gilbarg-trudinger:page-0089}
\dist(x,\mathcal N)\le K\|Sx\| \qquad \text{for all }x\in\V.
\end{equation}
\textit{Proof.} Suppose the result is not true. Then there exists a sequence $\{x_n\}\subset\V$ satisfying $\|Sx_n\|=1$ and $d_n=\dist(x_n,\mathcal N)\to\infty$. Choose a sequence $\{y_n\}\subset\mathcal N$ such that $d_n\le \|x_n-y_n\|\le 2d_n$. Then if
\[
z_n=\frac{x_n-y_n}{\|x_n-y_n\|}
\]
\textit{we have $\|z_n\|=1$, and $\|Sz_n\|\le d_n^{-1}\to0$ so that the sequence $\{Sz_n\}$ converges to $0$. But since $T$ is compact, by passing to a subsequence if necessary, we may assume that the sequence $\{Tz_n\}$ converges to an element $y_0\in\V$. Since $z_n=(S+T)z_n$, we then also have $\{z_n\}$ converging to $y_0$ and consequently $y_0\in\mathcal N$. However this leads to a contradiction as}
\[
\begin{aligned}
\dist(z_n,\mathcal N)
&=\inf_{y\in\mathcal N}\|z_n-y\|\\
&=\|x_n-y_n\|^{-1}\inf_{y\in\mathcal N}\|x_n-y_n-\|x_n-y_n\|\,y\|\\
&=\|x_n-y_n\|^{-1}\dist(x_n,\mathcal N)\ge \frac12.
\end{aligned}
\]

\item[(2)] \textit{Let $\mathcal R=S(\mathcal V)$ be the range of $S$. Then $\mathcal R$ is a closed subspace of $\V$.}

\textit{Proof.} Let $\{x_n\}$ be a sequence in $\V$ whose image $\{Sx_n\}$ converges to an element $y\in\V$. To show that $\mathcal R$ is closed we must show that $y=Sx$ for some element $x\in\V$. By our previous result the sequence $\{d_n\}$ where $d_n=\dist(x_n,\mathcal N)$ is bounded. Choosing $y_n\in\mathcal N$ as before and writing $w_n=x_n-y_n$, we consequently have that the sequence $\{w_n\}$ is bounded while the sequence $\{Sw_n\}$ converges to $y$. Since $T$ is compact, by passing to a subsequence if necessary we may assume that the sequence $\{Tw_n\}$ converges to an element $w_0\in\V$. Hence the sequence $\{w_n\}$ itself converges to $y+w_0$ and by the continuity of $S$, we have $S(y+w_0)=y$. Consequently $\mathcal R$ is closed. $\square$

\item[(3)] \textit{If $\mathcal N=\{0\}$, then $\mathcal R=\V$. That is, if case (i) of Theorem 5.3 does not hold, then case (ii) is true.}

\textit{Proof.} By our previous result the sets $\mathcal R_j$, defined by $\mathcal R_j=S^j(\mathcal V)$, $j=1,2,\dots$ form a non-increasing sequence of closed subspaces of $\V$. Suppose that no two of these spaces coincide. Then each is a proper subspace of its predecessor. Hence by Lemma 5.4, there exists a sequence $\{y_n\}\subset \V$ such that $y_n\in \mathcal R_n$, $\|y_n\|=1$ and $\mathrm{dist}\,(y_n,\mathcal R_{n+1})\geq \frac12$. Thus if $n>m$,
\[
T y_m - T y_n = y_m + (-y_n - S y_m + S y_n)
= y_m - y \quad \text{for some } y\in \mathcal R_{m+1}.
\]
Hence $\|T y_m - T y_n\|\geq \frac12$ contrary to the compactness of $T$. Consequently there exists an integer $k$ such that $\mathcal R_j=\mathcal R_k$ for all $j\geq k$. Up to this point we have not used the condition: $\mathcal N=\{0\}$. Now let $y$ be an arbitrary element of $\V$. Then $S^k y\in \mathcal R_k=\mathcal R_{k+1}$ and so $S^k y=S^{k+1}x$ for some $x\in \V$. Therefore $S^k(y-Sx)=0$ and so $y=Sx$ since $S^{-k}(0)=S^{-1}(0)=0$. Consequently $\mathcal R=\mathcal R_j=\V$ for all $j$. $\square$

\item[(4)] \textit{If $\mathcal R=\V$, then $\mathcal N=\{0\}$. Consequently either case (i) or case (ii) holds.}

\textit{Proof.} This time we define a non-decreasing sequence of closed subspaces $\{\mathcal N_j\}$ by setting $\mathcal N_j=S^{-j}(0)$. The closure of $\mathcal N_j$ follows from the continuity of $S$. By employing an analogous argument based on Lemma 5.4 to that used in step (3), we obtain that $\mathcal N_j=\mathcal N_l$ for all $j\geq \text{some integer } l$. Then if $\mathcal R=\V$, any element $y\in \mathcal N_l$ satisfies $y=S^l x$ for some $x\in \V$. Consequently $S^{2l}x=0$ so that $x\in \mathcal N_{2l}=\mathcal N_l$ whence $y=S^l x=0$. Step (4) is thus proved. $\square$
\end{enumerate}

The boundedness of the operator $S^{-1}=(I-T)^{-1}$ in case (ii) follows from step (1) with $\mathcal N=\{0\}$. Note that a slight simplification could be achieved by taking $\mathcal N=\{0\}$ at the outset in steps (1) and (2) and that step (4) is independent of the previous steps. Theorem 5.3 is thus completely proved. $\square$
\end{proof}

\begin{definition}\label{gt.def5.3.eigenvalue}\dcref{eigenvalue}\group{chapter-5}\source{gilbarg-trudinger:page-0090}
A number $\lambda$ is called an eigenvalue of $T$ if there exists a non-zero element $x$ in $\V$ (called an eigenvector) satisfying $Tx=\lambda x$. The dimension of the null space of the operator $S_\lambda=\lambda I-T$ is called the multiplicity of $\lambda$. If $\lambda\neq 0$, $\lambda\in \mathbb{R}$ is not an eigenvalue of $T$, it follows from Theorem 5.3 that the resolvent operator $R_\lambda=(\lambda I-T)^{-1}$ is a well defined, bounded linear mapping of $\V$ onto itself.
\end{definition}

\begin{theorem}\label{gt.thm5.5}\dcref{Theorem 5.5}\group{chapter-5}\source{gilbarg-trudinger:page-0090,page-0091}
A compact linear mapping $T$ of a normed linear space into itself possesses a countable set of eigenvalues having no limit points except possibly $\lambda=0$. Each non-zero eigenvalue has finite multiplicity.
\end{theorem}

\begin{proof}\source{gilbarg-trudinger:page-0090,page-0091}
Suppose that there exists a sequence $\{\lambda_n\}$ of not necessarily distinct eigenvalues and a sequence of corresponding linearly independent eigenvectors $\{x_n\}$ satisfying $\lambda_n\to \lambda\neq 0$. Let $\mathcal M_n$ be the closed subspace spanned by $\{x_1,\ldots,x_n\}$. By Lemma 5.4, there exists a sequence $\{y_n\}$ such that $y_n\in\mathcal M_n$, $\|y_n\|=1$ and $\operatorname{dist}(y_n,\mathcal M_{n-1})\geq \frac12$, $(n=2,3\ldots)$. If $n>m$, we have
\[
\lambda_n^{-1}Ty_n-\lambda_m^{-1}Ty_m
=y_n+(-y_m-\lambda_n^{-1}S_{\lambda_n}y_n+\lambda_m^{-1}S_{\lambda_m}y_m)
\]
\[
=y_n-z \qquad \text{where } z\in\mathcal M_{n-1}.
\]
For, if $y_n=\sum_{j=1}^{n}\beta_jx_j$ then $y_n-\lambda_n^{-1}Ty_n=\sum_{j=1}^{n}\beta_j(1-\lambda_n^{-1}\lambda_j)x_j\in\mathcal M_{n-1}$ and similarly $S_{\lambda_m}y_m\in\mathcal M_m$. Therefore we have
\[
\bigl\|\lambda_n^{-1}Ty_n-\lambda_m^{-1}Ty_m\bigr\|\geq \frac12
\]
which contradicts the compactness of $T$ combined with the hypothesis $\lambda_n\to\lambda\neq 0$.
Hence our initial supposition is false and this implies the validity of the theorem. $\square$
\end{proof}

\section*{5.4. \ Dual Spaces and Adjoints}
\source{gilbarg-trudinger:page-0091,page-0092}

\begin{definition}\label{gt.def5.4.dual-space}\dcref{dual space}\group{chapter-5}\source{gilbarg-trudinger:page-0091}
Let $\V$ be a normed linear space. A functional on $\V$ is a mapping from $\V$ into $\R$. The space of all bounded linear functionals on $\V$ is called the dual space of $\V$ and is denoted by $\V^*$. It can be shown easily that $\V^*$ is a Banach space under the norm
\begin{equation}\label{gt.eq5.5}\dcref{(5.5)}\group{chapter-5}\source{gilbarg-trudinger:page-0091}
\|f\|_{\V^*}=\sup_{x\neq 0}\frac{|f(x)|}{\|x\|}
\end{equation}
\end{definition}

\begin{example}\source{gilbarg-trudinger:page-0091}
The dual space of $\R^n$ is isomorphic to $\R^n$ itself.
\end{example}

\begin{definition}\label{gt.def5.4.reflexive}\dcref{reflexive}\group{chapter-5}\source{gilbarg-trudinger:page-0091}
The dual space of $\V^*$, denoted $\V^{**}$, is called the second dual of $\V$. Clearly the mapping $J:\V\to\V^{**}$ given by $Jx(f)=f(x)$ for $f\in\V^*$ is a norm preserving, linear, one-to-one mapping of $\V$ into $\V^{**}$. If $J\V=\V^{**}$, then we call $\V$ reflexive.
\end{definition}

\begin{definition}\label{gt.def5.4.adjoint}\dcref{adjoint}\group{chapter-5}\source{gilbarg-trudinger:page-0091}
Let $T$ be a bounded linear mapping between two Banach spaces $\B_1$ and $\B_2$. The adjoint of $T$, denoted $T^*$, is a bounded linear mapping between $\B_2^*$ and $\B_1^*$ defined by
\begin{equation}\label{gt.eq5.6}\dcref{(5.6)}\group{chapter-5}\source{gilbarg-trudinger:page-0091}
(T^*g)(x)=g(Tx) \qquad \text{for } g\in\B_2^*,\ x\in\B_1.
\end{equation}
\end{definition}

\begin{remark}\source{gilbarg-trudinger:page-0092}
Letting $\mathcal N$, $\mathcal R$, $\mathcal N^*$, $\mathcal R^*$ denote the null spaces and ranges of $T$, $T^*$ respectively, the following relations hold provided $\mathcal R$ is closed,
\[
\mathcal R={\Nstar}^{\perp}=\{\,y\in \B_2\mid g(y)=0 \quad \text{for all } g\in \Nstar\,\},
\]
\[
\mathcal R^*=\mathcal N^\perp=\{\,f\in \B_1^*\mid f(x)=0 \quad \text{for all } x\in \mathcal N\,\}.
\]
Also the compactness of $T$ implies the compactness of $T^*$. These two results are proved for example in [YO].
\end{remark}

\section*{5.5. \ Hilbert Spaces}
\source{gilbarg-trudinger:page-0092}

\begin{definition}\label{gt.def5.5.inner-product}\dcref{inner product}\group{chapter-5}\source{gilbarg-trudinger:page-0092}
A scalar (or inner) product on a linear space $\V$ is a mapping $q:\V\times\V\to \mathbb{R}$ (henceforth we write $q(x, y)=(x, y)$ or $(x, y)_{\V}$, $x, y\in\V$) satisfying
\begin{enumerate}
\item[(i)] $(x, y)=(y, x)$ for all $x, y\in\V$,
\item[(ii)] $(\lambda_1x_1+\lambda_2x_2, y)=\lambda_1(x_1, y)+\lambda_2(x_2, y)$ for all $\lambda_1, \lambda_2\in\mathbb{R}$, $x_1, x_2, y\in\V$,
\item[(iii)] $(x, x)>0$ for all $x\neq 0$, $x\in\V$.
\end{enumerate}
A linear space $\V$ equipped with an inner product is called an inner product space or a pre-Hilbert space. Writing $\|x\|=(x, x)^{1/2}$ for $x\in\V$, we have the following inequalities:
\end{definition}

\begin{equation}\label{gt.eq5.7}\dcref{(5.7)}\group{chapter-5}\source{gilbarg-trudinger:page-0092}
|(x, y)|\leq \|x\|\,\|y\|
\end{equation}
\begin{equation}\label{gt.eq5.8}\dcref{(5.8)}\group{chapter-5}\source{gilbarg-trudinger:page-0092}
\|x+y\|\leq \|x\|+\|y\|
\end{equation}
\begin{equation}\label{gt.eq5.9}\dcref{(5.9)}\group{chapter-5}\source{gilbarg-trudinger:page-0092}
\|x+y\|^2+\|x-y\|^2=2(\|x\|^2+\|y\|^2).
\end{equation}

\begin{definition}\label{gt.def5.5.hilbert-space}\dcref{Hilbert space}\group{chapter-5}\source{gilbarg-trudinger:page-0092}
In particular an inner product space $\V$ is a normed linear space. A Hilbert space is defined to be a complete inner product space.
\end{definition}

\begin{example}\source{gilbarg-trudinger:page-0093}
Euclidean space $\mathbb{R}^n$ is a Hilbert space under the inner product
\[
(x,y)=\sum x_i y_i,\qquad x=(x_1,\dots,x_n),\qquad y=(y_1,\dots,y_n).
\]
The Sobolev spaces $W^{k,2}(\Omega)$; (see Chapter 7).
\end{example}

\section*{5.6. \ The Projection Theorem}
\source{gilbarg-trudinger:page-0093,page-0094}

\begin{definition}\label{gt.def5.6.orthogonal}\dcref{orthogonal}\group{chapter-5}\source{gilbarg-trudinger:page-0093}
Two elements $x$ and $y$ in an inner product space are called orthogonal (or perpendicular) if $(x,y)=0$. Given a subset $\M$ of an inner product space we denote by $\M^\perp$ the set of elements orthogonal to every element of $\M$.
\end{definition}

\begin{theorem}\label{gt.thm5.6}\dcref{Theorem 5.6}\group{chapter-5}\source{gilbarg-trudinger:page-0093}
Let $\M$ be a closed subspace of a Hilbert space $\Hspace$. Then for every $x\in\Hspace$ we have $x=y+z$ where $y\in\M$ and $z\in\M^\perp$.
\end{theorem}

\begin{proof}\source{gilbarg-trudinger:page-0093}
If $x\in\M$, we set $y=x$, $z=0$. Hence we may assume $\M\neq \Hspace$ and $x\notin\M$. Define
\[
d=\dist(x,\M)=\inf_{y\in\M}\|x-y\|>0
\]
and let $\{y_n\}\subset\M$ be a minimizing sequence, that is $\|x-y_n\|\to d$. Using the parallelogram law we obtain
\[
4\left\|x-\frac{1}{2}(y_n+y_m)\right\|^2+\|y_n-y_m\|^2=2\left(\|x-y_n\|^2+\|x-y_m\|^2\right)
\]
so that, since $\frac{1}{2}(y_n+y_m)\in\M$, also we have $\|y_n-y_m\|\to 0$ as $m,n\to\infty$; that is the sequence $\{y_n\}$ converges since $\Hspace$ is complete. Also, since $\M$ is closed, $y=\lim y_n\in\M$ and $\|x-y\|=d$.

Now write $x=y+z$ where $z=x-y$. To complete the proof we must show $z\in\M^\perp$. For any $y'\in\M$ and $\alpha\in\mathbb{R}$ we have $y+\alpha y'\in\M$ and so
\[
d^2\le \|x-y-\alpha y'\|^2=(z-\alpha y',\,z-\alpha y')
\]
\[
=\|z\|^2-2\alpha (y',z)+\alpha^2\|y'\|^2.
\]
Therefore, since $\|z\|=d$, we obtain for all $\alpha>0$
\[
|(y',z)|\le \frac{\alpha}{2}\|y'\|^2
\]
so that $(y',z)=0$ for all $y'\in\M$. Hence $z\in\M^\perp$. $\square$
\end{proof}

\begin{remark}\source{gilbarg-trudinger:page-0094}
The element $y$ is called the orthogonal projection of $x$ on $\M$. Theorem 5.6 also shows that any closed proper subspace of $\Hspace$ is orthogonal to some element of $\Hspace$.
\end{remark}

\section*{5.7. \ The Riesz Representation Theorem}
\source{gilbarg-trudinger:page-0094}

\begin{theorem}\label{gt.thm5.7}\dcref{Theorem 5.7}\group{chapter-5}\source{gilbarg-trudinger:page-0094}
For every bounded linear functional $F$ on a Hilbert space $\Hspace$, there is a uniquely determined element $f\in \Hspace$ such that $F(x)=(x,f)$ for all $x\in \Hspace$ and $\|F\|=\|f\|$.
\end{theorem}

\begin{proof}\source{gilbarg-trudinger:page-0094}
Let $\mathcal N=\{x\mid F(x)=0\}$ be the null space of $F$. If $\mathcal N=\Hspace$, the result is proved by taking $f=0$. Otherwise, since $\mathcal N$ is a closed subspace of $\Hspace$, there exists by Theorem 5.6 an element $z\neq 0$, $z\in \Hspace$ such that $(x,z)=0$ for all $x\in \mathcal N$. Hence $F(z)\neq 0$ and moreover for any $x\in \Hspace$,
\[
F\!\left(x-\frac{F(x)}{F(z)}\,z\right)=F(x)-\frac{F(x)}{F(z)}F(z)=0
\]
so that the element $x-\frac{F(x)}{F(z)}\,z\in \mathcal N$. This means that
\[
\left(x-\frac{F(x)}{F(z)}\,z,z\right)=0,
\]
that is, that
\[
(x,z)=\frac{F(x)}{F(z)}\|z\|^2
\]
and hence $F(x)=(f,x)$ where $f=zF(z)/\|z\|^2$. The uniqueness of $f$ is easily proved and is left to the reader. To show that $\|F\|=\|f\|$, we have first, by the Schwarz inequality,
\[
\|F\|=\sup_{x\neq 0}\frac{|(x,f)|}{\|x\|}\leq \sup_{x\neq 0}\frac{\|x\|\,\|f\|}{\|x\|}=\|f\|;
\]
and secondly,
\[
\|f\|^2=(f,f)=F(f)\leq \|F\|\,\|f\|,
\]
so that $\|f\|\leq \|F\|$, and hence $\|F\|=\|f\|$. $\square$
\end{proof}

\section*{5.8. \ The Lax-Milgram Theorem}
\source{gilbarg-trudinger:page-0095}

\begin{definition}\label{gt.def5.8.bilinear}\dcref{bounded, coercive bilinear form}\group{chapter-5}\source{gilbarg-trudinger:page-0095}
A bilinear form $\mathbf{B}$ on a Hilbert space $\Hilb$ is called bounded if there exists a constant $K$ such that
\begin{equation}\label{gt.eq5.10}\dcref{(5.10)}\group{chapter-5}\source{gilbarg-trudinger:page-0095}
|\mathbf{B}(x,y)|\le K\|x\|\,\|y\| \qquad \text{for all $x,y\in \Hilb$}
\end{equation}
and coercive if there exists a number $v>0$ such that
\begin{equation}\label{gt.eq5.11}\dcref{(5.11)}\group{chapter-5}\source{gilbarg-trudinger:page-0095}
\mathbf{B}(x,x)\ge v\|x\|^2 \qquad \text{for all $x\in \Hilb$.}
\end{equation}
\end{definition}

\begin{theorem}\label{gt.thm5.8}\dcref{Theorem 5.8}\group{chapter-5}\source{gilbarg-trudinger:page-0095}
Let $\mathbf{B}$ be a bounded, coercive bilinear form on a Hilbert space $\Hilb$. Then for every bounded linear functional $F\in \Hilb^*$, there exists a unique element $f\in \Hilb$ such that
\[
\mathbf{B}(x,f)=F(x)\qquad \text{for all $x\in \Hilb$.}
\]
\end{theorem}

\begin{proof}\source{gilbarg-trudinger:page-0095}
By virtue of Theorem 5.7, there exists a linear mapping $T:\Hilb\to\Hilb$ defined by $\mathbf{B}(x,f)=(x,Tf)$ for all $x\in \Hilb$. Furthermore $\|Tf\|\le K\|f\|$ by \eqref{gt.eq5.10} so that $T$ is bounded. By \eqref{gt.eq5.11} we obtain $v\|f\|^2\le \mathbf{B}(f,f)=(f,Tf)\le \|f\|\,\|Tf\|,$ so that
\[
v\|f\|\le \|Tf\|\le K\|f\| \qquad \text{for all } f\in \Hilb.
\]
This estimate implies that $T$ is one-to-one, has closed range (see Problem 5.3) and that $T^{-1}$ is bounded. Suppose that $T$ is not onto $\Hilb$. Then there exists an element $z\ne 0$ satisfying $(z,Tf)=0$ for all $f\in \Hilb$. Choosing $f=z$, we obtain $(z,Tz)=\mathbf{B}(z,z)=0$ implying $z=0$ by \eqref{gt.eq5.11}. Consequently $T^{-1}$ is a bounded linear mapping on $\Hilb$. We then have $F(x)=(x,g)=\mathbf{B}(x,T^{-1}g)$ for all $x\in \Hilb$ and some unique $g\in \Hilb$ and the result is proved with $f=T^{-1}g$. $\square$
\end{proof}

\section*{5.9. \ The Fredholm Alternative in Hilbert Spaces}
\source{gilbarg-trudinger:page-0095,page-0096,page-0097}

\begin{definition}\label{gt.def5.9.adjoint-hilbert}\dcref{adjoint in a Hilbert space}\group{chapter-5}\source{gilbarg-trudinger:page-0096}
If $T$ is a bounded linear operator in a Hilbert space $\Hspace$, its adjoint $T^*$ is also a bounded linear mapping in $\Hspace$ defined by
\begin{equation}\label{gt.eq5.12}\dcref{(5.12)}\group{chapter-5}\source{gilbarg-trudinger:page-0096}
(T^*y,x)=(y,Tx) \quad \text{for all } x,y \in \Hspace.
\end{equation}
Clearly $\|T^*\|=\|T\|$, where $\|T\|=\sup_{x\neq 0}\|Tx\|/\|x\|$.
\end{definition}

\begin{lemma}\label{gt.lem5.9}\dcref{Lemma 5.9}\group{chapter-5}\source{gilbarg-trudinger:page-0096}
If $T$ is compact, then $T^*$ is also compact.
\end{lemma}

\begin{proof}\source{gilbarg-trudinger:page-0096}
Let $\{x_n\}$ be a sequence in $\Hspace$ satisfying $\|x_n\|\le M$. Then
\[
\|T^*x_n\|^2=(T^*x_n,T^*x_n)=(x_n,TT^*x_n)
\le \|x_n\|\,\|TT^*x_n\|
\le M\|T\|\,\|T^*x_n\|,
\]
so that $\|T^*x_n\|\le M\|T\|$; that is, the sequence $\{T^*x_n\}$ is also bounded. Hence, since $T$ is compact, by passing to a subsequence if necessary, we may assume that the sequence $\{TT^*x_n\}$ converges. But then
\[
\|T^*(x_n-x_m)\|^2=(T^*(x_n-x_m),T^*(x_n-x_m))
=(x_n-x_m,TT^*(x_n-x_m))
\le 2M\|T\|\,\|T^*(x_n-x_m)\|\to 0
\]
as $m,n\to\infty$.
Since $\Hspace$ is complete, the sequence $\{T^*x_n\}$ is convergent and hence $T^*$ is compact. $\square$
\end{proof}

\begin{lemma}\label{gt.lem5.10}\dcref{Lemma 5.10}\group{chapter-5}\source{gilbarg-trudinger:page-0096}
The closure of the range of $T$ is the orthogonal complement of the null space of $T^*$.
\end{lemma}

\begin{proof}\source{gilbarg-trudinger:page-0096}
Let $\mathcal R$ be the range of $T$, $\Nstar$ the null space of $T^*$. If $y=Tx$, we have
\[
(y,f)=(Tx,f)=(x,T^*f)=0 \quad \text{for all } f\in \Nstar
\]
so that $\mathcal R\subset {\Nstar}^{\perp}$, and since ${\Nstar}^{\perp}$ is closed, $\overline{\mathcal R}\subset {\Nstar}^{\perp}$. Now suppose that $y\in \Hspace$. By the projection theorem, Theorem 5.7, $y=y_1+y_2$ where $y_1\in \overline{\mathcal R},\, y_2\in \overline{\mathcal R}^{\,\perp}$. Consequently $(y_2,Tx)=(T^*y_2,x)=0$ for all $x\in \Hspace$, so that $y_2\in \Nstar$. Therefore $(y_2,y)=(y_2,y_1)+\|y_2\|^2=\|y_2\|^2$ and hence $y\in {\Nstar}^{\perp}$. $\square$
\end{proof}

\begin{theorem}\label{gt.thm5.11}\dcref{Theorem 5.11}\group{chapter-5}\source{gilbarg-trudinger:page-0096,page-0097}
Let $\Hspace$ be a Hilbert space and $T$ a compact mapping of $\Hspace$ into itself. Then there exists a countable set $\Lambda \subset \mathbb{R}$ having no limit points except possibly $\lambda=0$, such that if $\lambda\neq 0$, $\lambda\in \Lambda$, the equations
\begin{equation}\label{gt.eq5.13}\dcref{(5.13)}\group{chapter-5}\source{gilbarg-trudinger:page-0096,page-0097}
\lambda x-Tx=y,\ \lambda x-T^*x=y
\end{equation}
have uniquely determined solutions $x\in \Hspace$ for every $y\in \Hspace$, and the inverse mappings $(\lambda I-T)^{-1}$, $(\lambda I-T^*)^{-1}$ are bounded. If $\lambda\in\Lambda$, the null spaces of the mappings $\lambda I-T$, $\lambda I-T^*$ have positive finite dimension and the equations \eqref{gt.eq5.13} are solvable if and only if $y$ is orthogonal to the null space of $\lambda I-T^*$ in the first case and $\lambda I-T$ in the other.
\end{theorem}

\section*{5.10. \ Weak Compactness}
\source{gilbarg-trudinger:page-0097}

\begin{definition}\label{gt.def5.10.weak-convergence}\dcref{weak convergence}\group{chapter-5}\source{gilbarg-trudinger:page-0097}
Let $\V$ be a normed linear space. A sequence $\{x_n\}$ converges weakly to an element $x\in \V$ if $f(x_n)\to f(x)$ for all $f$ in the dual space $\V^*$. By the Riesz representation theorem, Theorem 5.7, a sequence $\{x_n\}$ in a Hilbert space $\Hspace$ will converge weakly to $x\in \Hspace$ if $(x_n,y)\to(x,y)$ for all $y\in \Hspace$.
\end{definition}

\begin{theorem}\label{gt.thm5.12}\dcref{Theorem 5.12}\group{chapter-5}\source{gilbarg-trudinger:page-0097}
A bounded sequence in a Hilbert space contains a weakly convergent subsequence.
\end{theorem}

\begin{proof}\source{gilbarg-trudinger:page-0097}
Let us assume initially that $\Hspace$ is separable and suppose that the sequence $\{x_n\}\subset \Hspace$ satisfies $\|x_n\|\leq M$. Let $\{y_m\}$ be a dense subset of $\Hspace$. By the Cantor diagonal process we obtain a subsequence $\{x_{n_k}\}$ of our original sequence satisfying $(x_{n_k},y_m)\to \alpha_m\in \mathbb{R}$ as $k\to\infty$. The mapping $f:\{y_m\}\to \mathbb{R}$ defined by $f(y_m)=\alpha_m$ may consequently be extended to a bounded linear functional $f$ on $\Hspace$ and hence by the Riesz representation theorem, there exists an element $x\in \Hspace$ satisfying $(x_{n_k},y)\to f(y)=(x,y)$ as $k\to\infty$, for all $y\in \Hspace$. Hence the subsequence $\{x_{n_k}\}$ converges weakly to $x$.
To extend the result to an arbitrary Hilbert space $\Hspace$, we let $\Hspace_0$ be the closure of the linear hull of the sequence $\{x_n\}$. Then by our previous argument there exists a subsequence $\{x_{n_k}\}\subset \Hspace_0$ and an element $x\in \Hspace_0$ satisfying $(x_{n_k},y)\to (x,y)$ for all $y\in \Hspace_0$. But by Theorem 5.5, we have for arbitrary $y\in \Hspace$, $y=y_0+y_1$, where $y_0\in \Hspace_0$, $y_1\in \Hspace_0^\perp$. Hence $(x_{n_k},y)=(x_{n_k},y_0)\to (x,y_0)=(x,y)$ for all $y\in \Hspace$ so that $\{x_{n_k}\}$ converges weakly to $x$, as required. $\square$
\end{proof}

\begin{remark}\source{gilbarg-trudinger:page-0097}
The first part of the proof of Theorem 5.12 extends automatically to reflexive Banach spaces with separable dual spaces (see Problem 5.4). The result is true however for arbitrary reflexive Banach spaces (see [YO]).
\end{remark}

\begin{remark}\source{gilbarg-trudinger:page-0097}
The material in this chapter is standard and can be found in texts on functional analysis such as [DS], [EW] and [YO].
\end{remark}

\section*{Problems}
\source{gilbarg-trudinger:page-0098}

\begin{enumerate}
\item[\textbf{5.1.}] Prove that the H\"older spaces $C^{k,\alpha}(\overline{\Omega})$, introduced in Chapter 4, are Banach spaces under either of the equivalent norms (4.6) or (4.6)$'$.
\item[\textbf{5.2.}] Prove that the interior H\"older spaces $C^{k,\alpha}_{*}(\Omega)$ defined by
\[
C^{k,\alpha}_{*}(\Omega)=\left\{u \in C^{k,\alpha}(\Omega)\,|\,|u|^{*}_{k,\Omega}<\infty\right\}
\]
are Banach spaces under the interior norms given by (4.17).
\item[\textbf{5.3.}] Let $\mathcal{B}$ be a Banach space and $T$ be a bounded linear mapping of $\mathcal{B}$ into itself satisfying
\[
\|x\| \leq K \|Tx\| \qquad \text{for all } x \in \mathcal{B},
\]
for some $K \in \mathbb{R}$. Prove that the range of $T$ is closed.
\item[\textbf{5.4.}] Prove that a bounded sequence in a separable, reflexive Banach space contains a weakly convergent subsequence.
\end{enumerate}

\section*{Chapter 6}
\source{gilbarg-trudinger:page-0099}

\section*{Classical Solutions; the Schauder Approach}
\source{gilbarg-trudinger:page-0099}

This chapter develops a theory of second order linear elliptic equations that is essentially an extension of potential theory. It is based on the fundamental observation that equations with H\"older continuous coefficients can be treated locally as a perturbation of constant coefficient equations. From this fact Schauder [SC 4, 5] was able to construct a global theory, an extension of which is presented here. Basic to this approach are apriori estimates of solutions, extending those of potential theory to equations with H\"older continuous coefficients. These estimates provide compactness results that are essential for the existence and regularity theory, and since they apply to classical solutions under relatively weak hypotheses on the coefficients, they play an important part in the subsequent nonlinear theory.

Throughout this chapter we shall denote by $Lu=f$ the equation
\begin{equation}\label{gt.eq6.1}\dcref{(6.1)}\group{chapter-6}\source{gilbarg-trudinger:page-0099}
Lu=a^{ij}(x)D_{ij}u+b^i(x)D_i u+c(x)u=f(x),\qquad a^{ij}=a^{ji},
\end{equation}
where the coefficients and $f$ are defined in an open set $\Omega\subset \mathbb{R}^n$ and, unless otherwise stated, the operator $L$ is strictly elliptic; that is,
\begin{equation}\label{gt.eq6.2}\dcref{(6.2)}\group{chapter-6}\source{gilbarg-trudinger:page-0099}
a^{ij}(x)\xi_i\xi_j\geq \lambda|\xi|^2,\qquad \forall x\in\Omega,\ \xi\in \mathbb{R}^n,
\end{equation}
for some positive constant $\lambda$.

\medskip

\noindent\textit{Equations with constant coefficients.} Before treating equation \eqref{gt.eq6.1} with variable coefficients, we establish a necessary preliminary result that extends Theorems 4.8 and 4.12 from Poisson's equation to other elliptic equations with constant coefficients. We state these extensions in the following lemma, recalling the interior and partially interior norms defined in (4.17), (4.18) and (4.29). Here and throughout this chapter all H\"older exponents will be assumed to lie in $(0,1)$ unless otherwise stated.

\medskip

\begin{lemma}\label{gt.lem6.1}\dcref{Lemma 6.1}\group{chapter-6}\source{gilbarg-trudinger:page-0099}
In the equation
\begin{equation}\label{gt.eq6.3}\dcref{(6.3)}\group{chapter-6}\source{gilbarg-trudinger:page-0099}
L_0u=A^{ij}D_{ij}u=f(x),\qquad A^{ij}=A^{ji},
\end{equation}
\end{lemma}
